\setcounter{section}{8}
\phantomsection
\section*{4. Conclusion}
\addcontentsline{toc}{section}{4. Conclusion}

This dissertation studied the relationship between birth order and academic productivity among historical European scholars. Specifically, using a new dataset of 645 pairs of brothers, the analysis examined whether older brothers showed higher scholarly productivity than their younger siblings and whether these patterns varied across historical periods and institutional environments.

The analysis provides evidence of a slight birth-order difference in scholarly productivity. The estimates indicate a modest but consistent productivity advantage for older brothers across several indicators of scholarly visibility and academic output. Crucially, birth-order differences appear to operate primarily through the extensive margin of scholarly activity: older siblings are more likely to enter scholarly production, while, conditional on publishing, productivity differences between brothers become substantially smaller and statistically insignificant. The results also reveal substantial historical heterogeneity. Before the sixteenth century, productivity differences between brothers appear more closely associated with sibling age gaps, whereas after the sixteenth century, differences in institutional affiliation become more strongly associated with productivity divergence within families. Because the analysis relies on observed scholarly siblings rather than complete family genealogies, the results should be interpreted as evidence on birth-order differences among observed academic elites rather than as direct estimates of strict primogeniture effects. At the same time, the results should be interpreted in light of the selected nature of historical scholarly records and the limited availability of detailed family information for earlier periods.

These results contribute to the existing literature. First, they extend the birth-order literature to a historical population of European intellectual elites, providing evidence that within-family productivity differences associated with birth order were already observable among pre-industrial academic elites. Second, focus on elite formation and intergenerational transmission by showing how sibling hierarchy, age structure, and institutional environments influenced scholarly trajectories. Finally, the results indicate that institutional and historical contexts were more relevant than inheritance systems in determining productivity differences among scholarly siblings, mainly through differences in institutional affiliation. These findings are also consistent with recent historical evidence suggesting that birth-order effects are limited in the general population but may nevertheless emerge more clearly within selected elite populations such as scholars \citep{ClarkCummins2024}.

Future research might extend the analysis to larger sibling groups, allowing for direct comparisons between firstborn, middle-born, and later-born children (that is, both secondborn siblings as well as those born after them in such a group). It can also focus on studying the mechanisms underlying these patterns more directly, including institutional mobility, academic networks, and intergenerational transmission of scholarly capital. Extending the analysis to other elite and non-elite subgroups might clarify the extent to which birth-order dynamics shaped long-run human-capital formation and social mobility in historical societies.